\begin{spacing}{1.5}
\vspace{1cm}

\textit{This chapter presents the results obtained using the methodology and implementation described earlier, and interprets what they mean. Organize this chapter around your research objectives/questions from Chapter 1 so the reader can clearly see whether each objective was met. Replace all italicized instruction text in this chapter with your own writing, and delete this paragraph.}

\section{Overview}

\textit{Briefly describe what experiments/evaluations were performed and outline what each section of this chapter covers.}

\section{Results}

\textit{Present your raw/quantitative results using tables, charts, and figures. Use clear captions and reference each table/figure in the text (e.g. ``as shown in Table~\textbackslash ref\{tab:performance\}''). Example results table:}

% Example results table:
% \begin{table}[h]
% \centering
% \small
% \captionsetup{font=small}
% \caption{Performance metrics of different models.}
% \begin{tabular}{@{}lcccccc@{}}
% \toprule
% \textbf{Model} & \textbf{Accuracy (\%)} & \textbf{Loss} & \textbf{Precision (\%)} & \textbf{Recall (\%)} & \textbf{F1 Score (\%)} \\ \midrule
% Model A & 0.00 & 0.0000 & 0.0 & 0.0 & 0.0 \\ \midrule
% \end{tabular}
% \label{tab:performance}
% \end{table}

\section{Detailed Analysis}

\textit{Go beyond the raw numbers: compare results across models/configurations/conditions, explain trends, and relate them back to your objectives and the literature reviewed in Chapter 2. Add subsections for different aspects of the analysis as needed.}

\subsection{Discussion}

\textit{Interpret the results: what do they mean for your research problem? Were the objectives met? Discuss any unexpected findings, limitations of the results, and how they compare with existing work.}

\section{Summary}

\textit{Summarize the key findings of this chapter and how they set up the conclusions in Chapter 6.}

\end{spacing}
